% !TEX root = ../notes.tex

\documentclass[../notes.tex]{subfiles}

\begin{document}

\section{March 2}
Here we go.

\subsection{More on the Higgs Field}
We continue with $G=\op{GL}_n$. It is remarkable that the genus of $C_b$ is (generically) $n^2(g-1)+1$, which agrees with $\dim\op{Bun}_G^0(X)$. We are going to attempt to embed $\op{Bun}_G^0(X)$ into the Jacobian of $C_b$.
\begin{remark}
	Suppose that $b\in\mc B$ is of the form $p(E,\varphi)$ and so that $C_b$ is smooth and irreducible. Then $C_b$ admits a line bundle $\mc L$ which is an eigenbundle for $\varphi$. Indeed, for some $\lambda\in C$ lying over $x\in X$, simply let $\mc L_\lambda$ be the eigenspace of $\varphi$ acting on the space $E_x$; generically, $\varphi$ has distinct eigenvalues, so we do in fact produce a well-defined line bundle.
\end{remark}
\begin{remark}
	Note that $E$ and $\varphi$ can both be reconstructed from the data of $C$ and $\mc L$: one has $\mc E=\pi_*\mc L$, where $\pi\colon C\to X$ is the canonical projection (indeed, this simply glues together the eigenspaces of $\varphi$ from $\mc L$ together into $\mc E$). Further, $\pi_*\mc L$ admits a canonical Higgs field $\pi_*\mc L\to\pi_*\mc L\otimes K_X$, where each summand of $\pi_*\mc L$ receives the corresponding eigenvalue; this Higgs field goes back to $\varphi$.
\end{remark}
We are now ready to complete the proof of \Cref{thm:hitchin-general} in this case.
\begin{proof}
	To start, we restrict $p$ to a subset of $T^*\op{Bun}_G^0(X)$ of a fixed degree $\ell$ so that it is connected. We see that Higgs fields $(E,\varphi)$ are in bijection with our pairs $(C,\mc L)$. One can calculate that $\deg\mc L$ is some fixed degree $d$, so $p^{-1}(\{b\})$ embeds in $\op{Pic}^dC$, which is isomorphic to $\op{Jac}C$. It follows that
	\[\dim p^{-1}(\{b\})\le n^2(g-1)+1.\]
	However, $p$ is of relative dimension at least $n^2(g-1)+1$ by calculating the difference of the dimensions of the codomain and domain. We conclude that all the fibers of $p$ must be the same dimension, and $p$ must be dominant, completing our proof.
\end{proof}
\begin{remark}
	For $G=\op{SL}_n$, it turns out that $p^{-1}(\{b\})$ embeds into the kernel of $\op{Jac}C\to\op{Jac}X$, which is some kind of Prym variety. For other classical groups, a similar argument works, but the story is more complicated in full generality.
\end{remark}

\subsection{Parabolic Structures}
We would like to say something for curves $X$ of smaller genus because curves of high genus are rather complicated to do calculations with.
\begin{definition}[structure]
	Choose reductive groups $H\subseteq G$, and let $\mc E$ be a $G$-bundle on $X$. Then an \textit{$H$-structure} on $\mc E$ at some fixed $x\in X$ is an $H$-orbit in the fiber $\mc E_x$.
\end{definition}
\begin{notation}
	Choose several parabolic subgroups $P_1,\ldots,P_N$ and points $t_1,\ldots,t_N\in X$. Then we let
	\[\op{Bun}_G(X,t,P)\]
	be the moduli stack of $G$-bundles on $X$ with parabolic $P_i$-structure at $t_i$ for each $i$.
\end{notation}
\begin{example}
	For $G=\op{GL}_N$ with $P$ given by the block upper-triangular matrices with sizes $N_1+\cdots+N_s=N$, then a $P$-structure on a vector bundle $\mc E$ of rank $N$ means that we are fixing a partial flag (whose quotient dimensions are given by the $N_\bullet$s) at a point.
\end{example}
\begin{remark}
	There is a natural map $\op{Bun}_G(X,t,P)\to\op{Bun}_G(X)$ which simply forgets the parabolic structures. Thus, this is a fibration whose fibers are $G/P_1\times\cdots\times G/P_N$. It follows that
	\[\dim\op{Bun}_G(X,t,P)=\dim\op{Bun}_G(X)+\sum_i(\dim G-\dim P_i).\]
\end{remark}
To retell our story of Hitchin systems, our next goal is to introduce some open substack of ``stable'' bundles with minimal stabilizers (namely, it would be nice to have stabilizer $Z(G)$). This is doable for large $N$, even for our curves of low genus.
\begin{example}
	Take $X=\PP^1$, and work with trivial $G$-bundles with some parabolic structures. It follows that the stack of these bundles is
	\[\Bigg(\prod_{i=1}^NG/P_i\Bigg)\mathbin{\bigg/}G.\]
	For example, $G=\op{PGL}_2$ and $P$ is the Borel subgroup, then we take $N\ge3$, and we are looking at the quotient $\left(\PP^1\right)^N/\op{PGL}_2$. In the open locus where the last three points of $\PP^1$ are distinct, there is a unique way to move them to $\{0,1,\infty\}$ with the $\op{PGL}_2$-action; thus, this locus is isomorphic to the variety $\left(\PP^1\right)^{N-3}$.
\end{example}
From now on, we will assume that we have a suitable definition of stable bundle, whose moduli we (unsurprisingly) denote $\op{Bun}_G^0(X,t,P)$.

\subsection{Parabolic Hitchin Systems}
Let's now attempt to define our Hitchin systems.
\begin{remark}
	For a given $\mc E$, the tangent space
	\[T_{\mc E}^*\op{Bun}_G^0(X,t,P)\]
	should consist of those Higgs fields $\varphi\colon\mc E\to\mc E\otimes K_X$ which have at worst simple poles at the $t_i$s and with residues $\alpha$ that strictly preserve the flag $\{V_\bullet\}$ at $t_i$, meaning that $\alpha V_i\subseteq V_{i-1}$.
\end{remark}
\begin{definition}[Hitchin map]
	The \textit{Hitchin map} is now defined exactly as before as
	\[p(E,\varphi)\coloneqq(Q_1(\varphi),\ldots,Q_r(\varphi)).\]
\end{definition}
\begin{example}
	We continue with the case where $G=\op{PGL}_2$ and the $P_i$s are all $B$. A parabolic structure at $t_i$ is an element $s_i\in\PP\mc E_{t_i}$. Thus, we are asking for our Higgs fields $\varphi\colon\mc E\to\mc E\otimes K_X$ (which must have vanishing trace to descend to $\op{PGL}_2$) to have simple poles at $t_i$ with residue $A_i$ such that $A_is_i=0$. Now, we know that
	\[p(E,\varphi)=\tr\land^2\varphi.\]
	Note that $\varphi$ is trace-zero and admits a kernel at each $t_i$, so it is nilpotent. It follows that $p(\mc E,\varphi)$ is $-\frac12\tr\varphi^2$, and staring at our poles shows that this lands in $\mathrm H^0(X;K_X^{\otimes2}\otimes\OO(t_1)\otimes\cdots\otimes\OO(t_N))$.
\end{example}
Let's continue with the above example, where $X=\PP^1$, and we will work with trivial bundles. (Note that all bundles are trivial up to some twist.) Suppose are asking for parabolic structures at the points $t_1,\ldots,t_N\in X$, which we may as well move away from $\infty$. Now, after doing a partial fraction decomposition, having the required poles means that $\varphi$ looks like
\[\varphi=\sum_{i=1}^N\frac{A_i}{z-t_i},\]
where $A_i$ are the residues; recall that they must satisfy $A_i^2=0$ to strictly preserve the flag while having vanishing trace. Furthermore, in order to be regular at $\infty$, a coordinate change ($z\mapsto1/z$) has no pole at $\infty$. It follows that $\sum_iA_i=0$.

Let's put some coordinates down. Write
\[A_i=\begin{bmatrix}
	a_i & b_i \\
	c_i & -d_i
\end{bmatrix}\]
for each $i$. To strictly preserve our flag, it must send some vector $(y_i,1)$ to zero. A short calculation shows that this matrix must actually take the form
\[A_i=p_i\begin{bmatrix}
	y_i & -y_i^2 \\
	1 & -y_i
\end{bmatrix}.\]
Thus, we see that $\frac12\tr\varphi^2$ is
\[\frac12\tr\sum_{i,j}\frac{A_iA_j}{(z-t_i)(z-t_j)}\,(dz)^2.\]
When $i=j$, the term goes away because $A_i^2=0$. Thus, we are left with
\[\frac12\tr\varphi^2=\sum_{i<j}\frac{\tr A_iA_j}{(z-t_i)(z-t_j)}\,(dz)^2.\]
Another calculation shows that each summand is $p_ip_j(y_i-y_j)^2$. Doing a partial fraction decomposition then shows
\[\frac12\tr\varphi^2=\sum_i\frac{G_i}{z-t_i}\,(dz)^2,\]
where
\[G_i\coloneqq\sum_{j\ne i}\frac{p_ip_j(y_i-y_j)^2}{t_j-t_i}.\]
These expressions $G_i$ are ``Garnier'' Hamiltonians. One can calculate that $\{G_i,G_j\}=0$, and there are $N$ of them, so it may look like we have produced an integrable system. However, these $G_i$s are not independent! Recall $\sum_iA_i=0$, so it follows that $\sum_ip_i=\sum_ip_iy_i=\sum_ip_iy_i^2=0$. (These are ``moment map'' conditions coming from $\mu=0$, where $\mu\colon T^*\CC^N\to\mf{sl}_2^*$ is the moment map.) On the other hand, $\sum_iG_i=\sum_it_iG_i=\sum_it_i^2G_i=0$.
\begin{remark}
	What is going on here is that $T^*\left(\PP^1\right)^{N-3}$ is being found inside $T^*\left(\PP^1\right)^N/\op{PGL}_2$. Thus, we cannot hope to have $N$ independent functions because they all come from a smaller space!
\end{remark}
Nonetheless, after passing to the quotient, we obtain an integrable system. Formally speaking, our ``phase space'' $T^*\left(\PP^1\right)^N/\op{PGL}_2$ should be given by pairs $(y,p)\in\CC^{2N}$ given by $\sum_ip_i=\sum_ip_iy_i=\sum_ip_iy_i^2$, and $\op{PGL}_2$ acts on the $y_\bullet$s by fractional linear transformations and on the $p_\bullet$s dually.

\subsection{Twisted Structures}
Suppose that all of our parabolic subgroups are a fixed Borel subgroup $B$. Using the double quotient argument from before, we find that
\[\op{Bun}_G(X,t,P)=G(R)\backslash G(\CC((z_1)))\times\cdots\times G(\CC((z_n)))/(I_1\times\cdots\times I_N),\]
where $R=\CC[X\setminus\{t_1,\ldots,t_N\}]$, and $I\subseteq G[[z]]$ is the subalgebra of elements which at $z=0$ evaluate to an element of $B$ at $z=0$; in particular, $G(\CC[[z]])/I$ is just $G/B$. Now, we can choose characters $\lambda_i\colon\op{Lie}I\to\CC$ coming from $\mf h^*$, and we let $\lambda=(\lambda_1,\ldots,\lambda_N)$ be the tuple of characters. There is now a twisted Hamiltonian reduction $\mu^{-1}(\lambda)/(G(R)\times I_1\times\cdots\times I_N)$, which is some twisted cotangent bundle of $\op{Bun}_G^0$.
\begin{remark}
	Let's say something about this twisted Hamiltonian reduction. For some group $H$ acting on $Y$, we have $T^*(Y/H)=T^*Y/H$ by usual Hamiltonian reduction. Then at a point $y\in Y$, one finds that $T_y^*(Y/H)$ is the kernel of $T_y^*Y\to\mf h^*$, where $\mf h=\op{Lie}H$. But we can instead take a pre-image by any choice of vector in $\mf h^*$, thereby ``twisting'' our reduction.
\end{remark}
\begin{example}
	We return to the case of $G=\op{PGL}_2$, where $X=\PP^1$. In this case, the cotangent space at $(\mc E,s_1,\ldots,s_N)$ consists of those Higgs bundles $\varphi\colon\mc E\to\mc E\otimes K_X$ with simple poles at $t_i$ with residues $A_i$ at $t_i$ for which $A_i|_{s_i}$ is exactly given by the character $\lambda_i$. Redoing the calculation of the previous subsection, one finds that $A_i$ satisfies $A_i^2=\lambda_i^2\id$ and so takes the form
	\[\begin{bmatrix}
		-\lambda_i+p_iy_i & 2\lambda_iy_i-p_iy_i^2 \\
		p_i & \lambda_i-p_iy_i
	\end{bmatrix}.\]
	One then can calculate $\varphi_i=\sum_i\frac{A_i}{z-z_i}$ has an explicit $\frac12\tr\varphi^2$ to receive Hamiltonians
	\[G_i\coloneqq\sum_{j\ne i}\frac{(y_i-y_j)^2p_ip_j-2(\lambda_ip_j-\lambda_jp_i)(y_i-y_j)-2\lambda_i\lambda_j}{t_j-t_i}.\]
	This gives an integrable system on the twisted cotangent bundle of $\left(\PP^1\right)^{N-3}$, where our coordinate is $\omega=\sum_idy_i\land dp_i$. One can calculate (as before) that $\sum_iG_i=0$, and both $\sum_it_iG_i$ and $\sum_it_i^2G_i$ do not depend on $\lambda$. Additionally, we have moment map conditions $\sum_ip_i=\sum_i(y_ip_i-\lambda_i)=\sum_i\left(y_i^2p_i-2\lambda_iy_i\right)=0$.
\end{example}
The moral is that we have a family of integrable systems living on some family of cotangent bundles.

\end{document}